\section{Introducción}

Los sistemas masa--resorte--amortiguador constituyen una aproximación fundamental para estudiar vibraciones mecánicas y respuestas transitorias. En este trabajo se emplea esta estructura para construir una analogía simplificada del movimiento vertical de una rueda: la masa equivalente se conecta a un chasis considerado fijo mediante un elemento elástico y un amortiguador viscoso, mientras que las irregularidades del terreno se representan mediante una fuerza externa variable en el tiempo.

El objetivo es obtener el modelo matemático mediante las ecuaciones de Lagrange, expresarlo en el dominio de Laplace y evaluar su comportamiento dinámico mediante un programa en Julia. Se considerarán dos experimentos: una respuesta natural desde una condición inicial no nula y una respuesta forzada desde el equilibrio y el reposo. Esta separación permitirá observar tanto las propiedades intrínsecas del sistema como su reacción ante una perturbación externa.

El modelo no pretende reproducir la dinámica completa de un vehículo. Su alcance se limita a un grado de libertad y omite, entre otros efectos, el movimiento del chasis, la rotación de la rueda, la elasticidad independiente del neumático y las no linealidades propias de una suspensión real.
